\chapter*{Conclusion}
\addcontentsline{toc}{chapter}{Conclusion}

This thesis presented a grammar-based greybox fuzzing system for automated vulnerability detection in SQLite, integrating structural grammar design, coverage-guided feedback, and adaptive rule selection into a unified pipeline. The system was built on top of the Nautilus 2.0 fuzzer \cite{aschermann2019nautilus} and extended with a multi-armed bandit algorithm for grammar rule weighting. A structural primitives grammar comprising 475 rules across three successive versions was developed through systematic analysis of confirmed CVEs, encoding the SQL constructs most likely to trigger memory safety violations. Experiments spanning 53 fuzzing campaigns across three SQLite versions (3.30.1, 3.31.1, and 3.32.0) --- each containing at least one confirmed CVE --- were conducted to evaluate whether bandit-driven adaptive rule selection offers an advantage over uniform sampling for CVE-class vulnerability discovery.

Four key findings emerged from the experimental evaluation. First, uniform sampling discovered approximately twice as many unique crash classes as bandit sampling (mean 8.6 versus 4.2, $p \approx 0.01$, Mann-Whitney U test), demonstrating that broad exploration of the grammar space outperforms early exploitation for triggering the diverse bug classes associated with known CVEs. Second, both policies located their first crash within one to two seconds of campaign start, but their discovery rates diverged over time: the bandit's per-minute crash rate decayed by 58\% across a campaign, compared to 51\% for uniform sampling, indicating that the bandit converged prematurely to a narrow subset of high-reward rules. Third, grammar structural coverage emerged as the primary bottleneck for CVE reachability: four of the six target CVEs were unreachable under earlier grammar versions and became reachable only after grammar v3.2 introduced the necessary SQL primitives, namely window functions and self-referential generated columns. Fourth, and most broadly, grammar evolution through structural expansion had a substantially larger effect on aggregate bug discovery than the choice of sampling policy, because the grammar defines the reachable search space while the policy merely navigates within it \cite{wang2019superion, aschermann2019nautilus}.

Several limitations bound these findings and should temper the generality of the conclusions. The sample size of five independent runs per experimental cell provides limited statistical power for per-version breakdown, making it difficult to attribute observed differences to specific version characteristics rather than sampling variance. All experiments were conducted on SQLite, and it is not known whether the relative performance of bandit versus uniform sampling holds for other database management systems with structurally different SQL dialects. The bandit's exponential moving average learning rate was fixed rather than systematically tuned, which may have contributed to its premature convergence and potentially understates the ceiling of bandit performance. Finally, the thesis focused exclusively on grammar rule selection as the adaptive decision and did not implement a full reinforcement learning agent capable of selecting among mutation operators such as random subtree replacement, recursive expansion, and splice \cite{sutton2018reinforcement}.

Three directions warrant further investigation as natural continuations of this work. The most direct extension is the integration of a Deep Q-Network (DQN) agent that learns to select not only grammar rules but also mutation operators based on joint coverage and crash-frequency signals --- this constitutes the planned Phase 3 of the RL-Nautilus project and would test whether learned mutation policies can recover the convergence problem observed in the bandit experiments. Second, extending the experimental scope to longer campaign durations of several hours to days, and to additional DBMS targets including MySQL and PostgreSQL, would assess how the relative merits of adaptive versus uniform sampling evolve when the search space is explored more thoroughly and under different grammatical structures. Third, exploring grammar composition techniques that automatically combine structural primitives extracted from multiple CVE analyses could partially automate the grammar design process that this thesis performed manually, potentially closing the loop between vulnerability disclosure and fuzzing grammar construction and making the pipeline more accessible to practitioners without deep SQL internals expertise.
